
\chapter*{Prologo}

Nel 2023, Geoffrey Hinton, considerato uno dei padri dell'intelligenza artificiale moderna, si è dimesso da Google con un avvertimento che ha fatto riflettere l'intera comunità scientifica: temeva che l'intelligenza artificiale potesse ereditare e amplificare i nostri pregiudizi, le nostre paure irrazionali, i nostri errori sistematici di giudizio\footnote{Metz, C. (2023). ``'The Godfather of A.I.' Leaves Google and Warns of Danger Ahead''. \textit{The New York Times}, 1 maggio 2023.}. La sua preoccupazione non riguardava tanto quello che le macchine possono fare, quanto quello che potrebbero imparare da noi: i nostri modi distorti di vedere il mondo, i nostri pregiudizi inconsci, le nostre scorciatoie mentali che ci portano sistematicamente fuori strada.

Questa preoccupazione tocca il cuore di un paradosso che definisce la condizione umana contemporanea. Abbiamo cervelli straordinari capaci di decodificare il genoma, esplorare l'universo, creare sinfonie che commuovono e teorie matematiche di eleganza sublime. Eppure questi stessi cervelli ci fanno comprare cose inutili quando abbiamo fame, ci convincono che siamo tutti guidatori sopra la media, ci fanno credere a teorie del complotto palesemente false, e ci tengono svegli fino alle tre del mattino a scrollare contenuti che nemmeno ricorderemo il giorno dopo. Come può la stessa mente che ha risolto i misteri dell'atomo cadere in trappole mentali così elementari?

La risposta sta in una verità scomoda che raramente affrontiamo: viviamo in due tempi simultaneamente. I nostri corpi e le nostre tecnologie esistono nel ventunesimo secolo, ma i nostri cervelli operano ancora con software cognitivo sviluppato nel Pleistocene. Questo disallineamento temporale - millenni di evoluzione biologica che si scontrano con decenni di rivoluzione tecnologica - non è semplicemente una curiosità accademica. È la fonte di frizioni cognitive che permeano ogni aspetto della nostra vita quotidiana, dalle decisioni finanziarie alle relazioni personali, dalla politica alla salute mentale.
\\ 


Quando sentiamo parlare di distorsioni mentali sistematiche o semplicemente \textbf{"bias cognitivi"}, tendiamo a immaginarli come semplici errori della mente. Li consideriamo come degli sbagli da correggere per riuscire finalmente a pensare in modo perfettamente logico.
Ma questa prospettiva manca completamente il punto. I bias cognitivi non sono bug nel sistema; sono il sistema stesso. Sono meccanismi mentali forgiati da milioni di anni di selezione naturale, strumenti cognitivi che hanno garantito la sopravvivenza dei nostri antenati in ambienti ostili e imprevedibili. Il problema non è che abbiamo questi meccanismi; il problema è che li stiamo usando in un mondo radicalmente diverso da quello per cui furono sviluppati.

Pensate a quando entrate in un supermercato moderno. Il vostro cervello si trova di fronte a una situazione senza precedenti evolutivi. I nostri antenati non hanno mai dovuto scegliere tra trecento varietà di cereali per la colazione o decidere quale delle cinquanta marche di dentifricio comprare. La massima diversità che potevano incontrare era limitata a quello che l'ambiente immediato offriva stagionalmente: alcune varietà di frutti, qualche tipo di radice, la carne di pochi animali. Il cervello umano si è evoluto per prendere decisioni rapide in contesti di scarsità relativa, dove le opzioni erano limitate e le conseguenze immediate. Ora questo stesso cervello deve navigare un'abbondanza paralizzante di scelte, dove ogni decisione richiede la valutazione di variabili astratte come rapporto qualità-prezzo, sostenibilità ambientale, implicazioni etiche della produzione.

Il risultato è quello che Barry Schwartz ha chiamato il paradosso della scelta\footnote{Schwartz, B. (2004). \textit{The Paradox of Choice: Why More Is Less}. New York: Ecco.}: l'abbondanza di opzioni, invece di renderci più liberi e felici, genera ansia, paralisi decisionale e insoddisfazione cronica. Non è debolezza o inadeguatezza personale; è il nostro hardware cognitivo del Pleistocene che cerca disperatamente di gestire un software del ventunesimo secolo per cui non è stato progettato.

Questo disallineamento diventa ancora più problematico quando consideriamo la natura delle sfide moderne. Il nostro cervello eccelle nel valutare minacce immediate e concrete: un movimento sospetto tra i cespugli, un odore che segnala cibo avariato, un'espressione facciale che suggerisce ostilità. Ma come può questo stesso sistema valutare minacce astratte e a lungo termine come il cambiamento climatico, che si manifesta su scale temporali e spaziali che trascendono completamente la nostra esperienza sensoriale diretta? 

Come può un cervello calibrato per relazioni faccia a faccia in gruppi di cinquanta persone gestire interazioni sociali mediate digitalmente con migliaia di sconosciuti?

La risposta è che non può, non senza significative distorsioni e errori sistematici. Quando un investitore esperto commette errori elementari durante una crisi di mercato, non stiamo osservando un fallimento individuale di razionalità. Stiamo osservando l'attivazione di meccanismi di panico ancestrali - gli stessi che salvavano vite quando il pericolo era un predatore - applicati inappropriatamente a fluttuazioni numeriche su uno schermo. Quando un medico brillante ignora sintomi che contraddicono la sua diagnosi iniziale, non è negligenza professionale; è il bias di conferma - lo stesso meccanismo che aiutava i nostri antenati a mantenere coerenza nelle loro conoscenze di sopravvivenza - che opera in un contesto dove la flessibilità mentale sarebbe più appropriata. 

Daniel Kahneman, il cui lavoro pionieristico sui bias cognitivi gli valse il Premio Nobel per l'economia, ha elegantemente descritto come la nostra mente operi attraverso due sistemi distinti\footnote{Kahneman, D. (2011). \textit{Thinking, Fast and Slow}. New York: Farrar, Straus and Giroux.}.  
Il Sistema 1 è veloce, automatico, intuitivo, emotivo. È quello che vi fa riconoscere istantaneamente un volto familiare, saltare indietro quando vedete qualcosa che sembra un serpente, sentire che qualcosa "non quadra" in una situazione anche se non sapete dire esattamente cosa. Il Sistema 2 è lento, deliberativo, logico, faticoso. È quello che usate per risolvere problemi matematici complessi, valutare argomenti filosofici, o pianificare un progetto a lungo termine.

La maggior parte della nostra vita mentale è governata dal Sistema 1, e per buone ragioni evolutive. In un mondo dove la velocità di reazione poteva fare la differenza tra vita e morte, un sistema che potesse valutare istantaneamente situazioni e generare risposte immediate era essenziale. Il problema è che questo sistema opera attraverso euristiche - scorciatoie mentali - che erano adaptive in contesti ancestrali ma possono essere profondamente fuorvianti nel mondo moderno. Il Sistema 2, che potrebbe correggere questi errori, è metabolicamente costoso e quindi pigro. Si attiva solo quando assolutamente necessario, e anche allora spesso si limita a razionalizzare le conclusioni già raggiunte dal Sistema 1 piuttosto che metterle genuinamente in discussione.

Questa architettura cognitiva a due velocità crea situazioni paradossali. Dopo decenni di studio sui bias cognitivi, Kahneman ammise candidamente di non essere diventato significativamente migliore nell'evitarli nella propria vita. La conoscenza dei bias, scoprì, non conferisce immunità da essi. È come conoscere perfettamente le leggi dell'ottica ma continuare a vedere le illusioni ottiche: la comprensione intellettuale non modifica la percezione immediata. Questo non è un fallimento personale di Kahneman; è una dimostrazione di quanto profondamente questi meccanismi siano radicati nella struttura stessa del nostro cervello.

Il biologo François Jacob descrisse l'evoluzione non come un ingegnere che progetta soluzioni ottimali, ma come un "bricoleur" - un tuttofare che lavora con i materiali disponibili, modificando e adattando strutture esistenti per nuovi scopi\footnote{Jacob, F. (1977). "Evolution and tinkering". \textit{Science}, 196(4295), 1161-1166.}. Il cervello umano è il prodotto di questo bricolage evolutivo durato milioni di anni. Non è un organo unitario progettato con coerenza, ma una stratificazione geologica di strutture che si sono evolute in epoche diverse per rispondere a pressioni selettive diverse. Il tronco encefalico che condividiamo con i rettili gestisce funzioni basali di sopravvivenza. Il sistema limbico che condividiamo con tutti i mammiferi processa emozioni e memoria. La neocorteccia, particolarmente sviluppata nei primati e massimamente negli umani, permette pensiero astratto e pianificazione.

Questi sistemi non sono semplicemente sovrapposti come strati di una torta; sono profondamente interconnessi in modi che creano sia la nostra straordinaria flessibilità cognitiva sia le nostre vulnerabilità sistematiche. Quando proviamo paura irrazionale di fronte a stimoli innocui - un ragno innocuo in un paese dove non esistono ragni velenosi, un ascensore perfettamente sicuro e regolarmente ispezionato, un gruppo di adolescenti che chiacchierano innocuamente - non è la nostra corteccia razionale che risponde per prima. È \textbf{l'amigdala}, una struttura antica del sistema limbico, che attiva risposte di allarme millisecondi prima che la ragione possa anche solo iniziare la sua analisi. Questa precedenza temporale dell'emozione sulla ragione non è un difetto di design; è una caratteristica che ha salvato innumerevoli vite quando la velocità di reazione alla minaccia era più importante dell'accuratezza della valutazione.

Ma c'è una dimensione ancora più profonda che spesso trascuriamo: la natura fondamentalmente sociale dei nostri bias. Il cervello umano non si è evoluto in isolamento ma nel contesto di gruppi sociali complessi. Molti dei nostri bias più pervasivi - il \textbf{bias di conferma} che ci fa cercare informazioni che confermano quello che già crediamo, \textbf{l'effetto gregge} che ci spinge a seguire la maggioranza, il \textbf{bias di autorità} che ci fa dare peso eccessivo alle opinioni di chi percepiamo come esperto - sono meccanismi che facilitavano la coesione e il coordinamento del gruppo in contesti ancestrali.

Consideriamo il bias di conferma più da vicino. Nel mondo moderno, questo bias può portare a polarizzazione politica estrema, bolle informative impenetrabili, e incapacità cronica di correggere false credenze anche di fronte a evidenze schiaccianti. Ma nel contesto di piccoli gruppi di cacciatori-raccoglitori, questo stesso meccanismo serviva funzioni vitali. Mantenere coerenza nelle credenze del gruppo su quali piante fossero commestibili, quali territori sicuri, quali alleanze affidabili, era letteralmente questione di vita o morte. Il dubbio costante e la rivalutazione continua delle conoscenze consolidate avrebbero paralizzato l'azione collettiva in situazioni dove la velocità e la coordinazione erano essenziali.

Oggi viviamo in società di milioni di individui, esposti a flussi informativi di complessità e volume che nessun cervello umano può realmente processare completamente. Il bias di conferma che una volta manteneva coesione in gruppi di cinquanta persone ora crea frammentazione in società di milioni. I social media amplificano questo effetto in modi che stiamo solo iniziando a comprendere. Gli algoritmi di raccomandazione, ottimizzati per massimizzare il coinvolgimento, hanno scoperto che il modo più efficace per tenere le persone incollate agli schermi è mostrar loro contenuti che confermano e amplificano le loro credenze preesistenti, non importa quanto distorte o false possano essere.

Questo ci porta a una delle sfide più inquietanti del nostro tempo: l'intersezione tra bias cognitivi umani e intelligenza artificiale. Quando creiamo sistemi di IA addestrati su dati generati da menti umane afflitte da bias, inevitabilmente codifichiamo questi bias nei sistemi stessi. Un Large Language Model\footnote{ChatGPT è un esempio di Large Language Model (LLM), un'intelligenza artificiale capace di comprendere e generare testo in linguaggio naturale. È stato addestrato su grandi quantità di testi da libri, articoli e siti web, imparando schemi linguistici per produrre risposte coerenti.} addestrato su miliardi di pagine di testo umano non impara solo la grammatica e il vocabolario; assorbe anche tutti i pregiudizi impliciti, gli stereotipi culturali, le distorsioni sistematiche presenti in quei testi. Se nei dati di training gli amministratori delegati sono prevalentemente descritti come uomini e le segretarie come donne, il modello "impara" queste associazioni e le riproduce nelle sue risposte, perpetuando e potenzialmente amplificando bias di genere che stiamo cercando di superare come società.

Ma c'è qualcosa di ancora più sottile e preoccupante. Questi sistemi di AI non sperimentano la \textbf{dissonanza cognitiva} ,  quel disagio psicologico che proviamo quando manteniamo credenze contraddittorie o quando il nostro comportamento contraddice i nostri valori. Quando fumate sapendo che fa male alla salute, o quando vi arrabbiate con i vostri figli dopo aver promesso a voi stessi di essere pazienti, quella sensazione sgradevole che provate è la dissonanza cognitiva in azione.

Per noi umani, questo disagio funziona come un sistema di allarme mentale. Ci segnala che qualcosa nel nostro sistema di credenze non è coerente e ci spinge a risolvere la contraddizione: modifichiamo il comportamento, cambiamo le nostre convinzioni, o troviamo giustificazioni per ridurre il conflitto. È un meccanismo evolutivo che ci orienta verso la coerenza interna, anche se imperfetta.

Un'intelligenza artificiale, invece, può produrre affermazioni completamente contraddittorie in risposte successive senza alcun problema. Può sostenere una posizione e il suo esatto opposto senza sperimentare alcun disagio interno, senza alcun impulso biologico verso la coerenza. Non esiste in essa quella spinta evolutiva che nei mammiferi ,  e particolarmente negli umani ,  genera il bisogno di risolvere le contraddizioni cognitive. Questa assenza di conflitto interno le permette di mantenere posizioni incoerenti indefinitamente, un comportamento che per noi sarebbe psicologicamente insostenibile.

È curioso come, proprio mentre iniziamo a comprendere i limiti della nostra percezione mentale attraverso gli LLM, abbiamo anche scoperto i limiti della nostra percezione fisica dell'universo. Nel 2015, per la prima volta nella storia, abbiamo rilevato le onde gravitazionali ,  increspature dello spazio-tempo teorizzate da Einstein un secolo prima\footnote{Il rilevamento delle onde gravitazionali da parte dell'interferometro LIGO valse il Premio Nobel per la Fisica nel 2017 a Rainer Weiss, Barry C. Barish e Kip S. Thorne. L'Italia partecipa a questa rivoluzione scientifica con VIRGO, uno dei principali rilevatori al mondo, situato in provincia di Pisa.}. Per millenni, l'umanità ha osservato l'universo solo attraverso la luce, come se fossimo creature che potevano solo vedere ma non sentire. Le onde gravitazionali ci hanno dato un nuovo senso, permettendoci di "ascoltare" eventi cosmici che la luce non può raccontarci: la fusione di buchi neri a miliardi di anni luce, le vibrazioni primordiali dell'universo nascente, fenomeni che sfuggono completamente alla nostra percezione elettromagnetica.
Il parallelo è sorprendente. Come le onde gravitazionali ci rivelano una dimensione nascosta dell'universo fisico ,  eventi invisibili ma che deformano la struttura stessa dello spazio-tempo ,  così i Large Language Models ci rivelano una dimensione nascosta della mente umana: i pattern invisibili del nostro pensiero collettivo, le distorsioni sistematiche che deformano la nostra percezione della realtà. Entrambe le scoperte ci mostrano quanto fossimo ciechi a intere dimensioni dell'esistenza. Le onde gravitazionali viaggiano attraverso il cosmo quasi indisturbate per miliardi di anni, portando informazioni da epoche e luoghi altrimenti inaccessibili. Allo stesso modo, i bias cognitivi viaggiano attraverso le generazioni umane, trasmessi culturalmente e ora codificati digitalmente negli LLM, portando con sé informazioni sul nostro passato evolutivo che altrimenti rimarrebbero sepolte nell'inconscio collettivo.

Eppure, nonostante queste sfide formidabili, o forse proprio a causa di esse, viviamo in un momento di opportunità straordinaria. Per la prima volta nella storia umana, abbiamo gli strumenti per osservare come funziona davvero il nostro cervello. Le risonanze magnetiche ci mostrano quali aree si accendono quando prendiamo decisioni sbagliate. I big data rivelano i pattern dei nostri errori collettivi. E ora, con lo sviluppo dei LLM, abbiamo qualcosa di ancora più interessante: uno specchio digitale che ci mostra come pensiamo davvero, non come crediamo di pensare.

Questi modelli di intelligenza artificiale, addestrati su miliardi di testi umani, hanno assorbito non solo il nostro linguaggio ma anche i nostri modi di ragionare ,  compresi tutti i nostri difetti. Studiarli è come guardare una radiografia del pensiero umano: vediamo le fratture, le asimmetrie, i punti deboli che normalmente ci sfuggono.

Questa nuova consapevolezza non ci rende immuni dai nostri errori. Daniel Kahneman, dopo una vita passata a studiare i bias cognitivi, ammetteva candidamente di caderci ancora. È come conoscere le regole della dieta perfetta e trovarsi comunque davanti al frigorifero a mezzanotte. La conoscenza aiuta, ma non basta. Quello che possiamo fare è imparare a riconoscere i momenti in cui siamo più vulnerabili, le situazioni in cui i nostri automatismi mentali ci portano fuori strada.

\vspace{1em}

Il libro che avete tra le mani è strutturato come un viaggio in quattro tappe:

\begin{itemize}
    \item \textbf{Parte I: La Mente Inaspettata.} 
    Inizieremo dalle basi, esplorando come l'evoluzione ha costruito il nostro cervello. Scopriremo perché siamo programmati per vedere pericoli dove non ce ne sono e opportunità dove non dovremmo vederle.

    \item \textbf{Parte II: La Mente Fuori Posto.} 
    Osserveremo questa mente antica alle prese con il mondo moderno. Dal supermercato alla borsa, dai social media agli ospedali, vedremo come i nostri istinti preistorici creano problemi molto contemporanei.

    \item \textbf{Parte III: Lo Specchio Digitale.} 
    Esploreremo come l'intelligenza artificiale sta imparando da noi ,  compresi i nostri pregiudizi. È uno specchio che non solo riflette chi siamo, ma amplifica i nostri difetti in modi che stiamo solo iniziando a capire.

    \item \textbf{Parte IV: La Bussola del Pensiero.} 
    Dall'analisi passeremo all'azione. Non ricette miracolose, ma strumenti pratici per navigare un mondo progettato per hackerare la nostra mente paleolitica.
\end{itemize}

Questo libro non vi renderà perfetti. La perfezione cognitiva non esiste e, se esistesse, probabilmente non sarebbe nemmeno desiderabile. I nostri "difetti" mentali sono anche ciò che ci rende creativi, intuitivi, capaci di decisioni rapide in situazioni complesse. Il problema non sono i bias in sé, ma non sapere quando ci stanno guidando.

Vi invito quindi a leggere queste pagine con curiosità più che con l'ansia di correggervi. I bias che incontrerete non sono errori da eliminare ma caratteristiche da comprendere. Sono presenti nel cervello del premio Nobel come in quello del vostro vicino di casa. Sono il prezzo che paghiamo per avere un cervello che funziona in tempo reale con risorse limitate, che deve prendere migliaia di decisioni al giorno senza potersi fermare a calcolare le probabilità di ogni scelta.

In un'epoca in cui algoritmi progettati per sfruttare le nostre debolezze cognitive competono per la nostra attenzione, in cui decisioni prese in un istante possono avere conseguenze globali, capire come funziona la nostra mente non è più un lusso intellettuale. È una competenza di sopravvivenza.

Questo libro è una mappa per orientarsi nel territorio accidentato della mente umana. Non una mappa perfetta ,  quella non esiste ,  ma una guida pratica per riconoscere quando stiamo per cadere in una trappola mentale, quando il nostro cervello ancestrale sta prendendo decisioni per noi, quando sarebbe meglio fermarsi e pensare invece di seguire l'istinto.

L'esplorazione che state per iniziare potrebbe essere scomoda. Scoprirete che molte delle vostre decisioni "razionali" non lo erano affatto. Che opinioni di cui andate fieri sono il prodotto di meccanismi mentali automatici. Che giudicate gli altri per gli stessi errori che commettete voi. Ma questa consapevolezza, per quanto sgradevole, è il primo passo verso una vita mentale più ricca e consapevole.

Siamo esseri straordinari: abbastanza intelligenti da capire di non essere sempre intelligenti, abbastanza razionali da riconoscere la nostra irrazionalità. È questo paradosso che ci rende umani. Non siamo macchine perfette e non lo saremo mai. Ma forse è proprio questa imperfezione, questa tensione continua tra quello che siamo e quello che vorremmo essere, che rende la vita interessante.

Il viaggio inizia ora. Preparatevi a incontrare il sabotatore che vive nella vostra testa, a scoprire perché continuate a fare cose che sapete essere sbagliate, a capire perché la vostra mente ,  questo capolavoro dell'evoluzione ,  sembra così spesso lavorare contro di voi. Ma preparatevi anche a scoprire che questi "difetti" sono in realtà funzionalità, che il sabotatore sta solo cercando di proteggervi con strumenti obsoleti, che la vostra mente imperfetta è anche quella che vi permette di amare, creare, sognare.

Non ne uscirete "guariti" ,  non c'è niente da cui guarire. Ma è come quando qualcuno vi fa notare che respirate: all'improvviso ne diventate consapevoli, e per un attimo potete scegliere come farlo. Questo libro vuole fare la stessa cosa con i vostri pensieri. Poi tornerete a pensare in automatico, certo. Ma in quei momenti cruciali ,  quando state per comprare qualcosa di inutile, per litigare per niente, per credere a una bugia ,  forse vi ricorderete di respirare.